% =====================================================================
%  27.2  Mo on subtrees -- the preorder flatten makes subtree(v) the
%  contiguous range [tin(v), tout(v)], so 27.1's engine runs verbatim;
%  plus the solved root-path batch CF 463E, where the window is the DFS
%  stack itself and the <=50 updates legalise one rebuild per block.
%  Exports: sec:c21-subtree  def:c21-flatten  thm:c21-subtree-range
%  lst:c21-flatten  prob:c21-caisa  lst:c21-caisa
%  Solved CF 463E "Caisa and Tree" (2100, L4), live-verified against the
%  CF problem page (name, root-1 strict-ancestor gcd query, type-2
%  updates capped at 50, n,q <= 100000, a_i <= 2*10^6; official sample
%  hand-replayed to -1 1 2 -1 1 BEFORE writing).
%  CF 383C (parity + difference BIT) is PROSE-ONLY here -- it is Ch.18's
%  own solved template (\chref{c14}); CF 375D's URL lives in Ch.17.
%  No URL is re-linked; urlsweep stays clean.
% =====================================================================
\section{Mo on Subtrees: the Preorder Flatten}\label{sec:c21-subtree}

\begin{previouslybox}[27.1 handed you the sweep; Ch.~5 handed you the array]
\secref{c21-moengine} built the engine: block-sort offline ranges, creep two
pointers, keep a statistic with \(O(1)\) \texttt{add}/\texttt{rm}, and every
range is answered in one \(\bigO{(n+q)\sqrt n}\) sweep.  What it never needed
was a tree -- so this section bolts the tree on, in the cheapest possible way.
Chapter~5 (\chref{c24}) already showed that one DFS fixes \(\tin(v),\tout(v)\),
and that the preorder array \(P\) with \(P[\tin(v)]=v\) turns \emph{every}
subtree into a contiguous segment.  A subtree query is therefore \emph{already}
an array range query: hand it to the 27.1 sweep unchanged.  And when the
statistic \emph{does} telescope (the sign-flip subtree update Ch.~18 solves as
CF 383C, \chref{c14}), you skip Mo entirely and spend a Fenwick difference
array instead.  The solved CF 463E below opens a \emph{third} door: the queried
set is not a subtree range but the current root path -- and a plain DFS carries
that window for free.  Knowing which world you live in is half the section.
\end{previouslybox}

\begin{topicintro}
Two lemmas do all the work.  \defref{c21-flatten} recalls the flatten: the DFS
entry/exit pair \([\tin(v),\tout(v)]\) is exactly the index span of
subtree\((v)\) inside preorder \(P\).  \thmref{c21-subtree-range} upgrades it:
running the \secref{c21-moengine} sweep on ranges \([\tin(v),\tout(v)]\) answers
\(q\) offline subtree statistics in \(\bigO{(n+q)\sqrt n}\) -- the engine does
not care that the array was baked from a tree.  The \emph{worked} running
example flattens \(R\) and reads subtree\((2)=P[2..5]\) off one line.  The
\solvedproblem{c21-caisa}{Caisa and Tree}{CF 463E}
  {https://codeforces.com/problemset/problem/463/E}{CF 2100 L4}
  {sieve of smallest prime factors; the DFS stack as the root-path window; a per-prime deepest-ancestor table with \(O(1)\) push/undo; one rebuild per block of the rare (\(\le 50\)) updates}
then opens the third world: "deepest strict ancestor sharing a prime factor
with \(v\)" has no inverse and no monotone lifting predicate -- but the queried
root paths are the recursion stack itself, so a plain DFS carries the window,
and the promise of at most \(50\) updates makes one rebuild per query block the
cheapest move on the board.

\textbf{Prerequisites.} \secref{c21-moengine} (the sweep, \texttt{add}/\texttt{rm});
\chref{c24} (\(\tin/\tout\), the flatten); \chref{c14} (Ch.~18, Fenwick trees --
where the invertible subtree-update template CF 383C is fully solved).

\textbf{After this section you can}
\begin{itemize}[nosep]
  \item flatten any rooted tree and name the range \([\tin(v),\tout(v)]\) of a subtree; \hfill\lv{3}
  \item run Mo on subtree ranges for a non-invertible statistic (distinct colours); \hfill\lv{4}
  \item carry a root-path statistic on the DFS stack with an \(O(1)\) push/undo table; \hfill\lv{4}
  \item solve CF 463E cold -- solved, traced to \(-1,1,2,-1,1\). \hfill\lv{4}
\end{itemize}
\end{topicintro}

% ---------------------------------------------------------------
\subsection{The flatten, and Mo rides it}
% ---------------------------------------------------------------

\begin{definition}{Preorder flatten and the subtree span}{c21-flatten}
Root the tree, run the left-to-right DFS, and timestamp \(\tin(v)\) = arrival
order (\(1\)..\(n\)) and \(\tout(v)\) = the last arrival inside \(v\)'s region.
Write \(P[\tin(v)]=v\).  Then
\[
  \subtree(v) \;=\; \{\,u : \tin(v)\le \tin(u)\le \tout(v)\,\},
\]
so \emph{aggregate over subtree\((v)\)} \(=\) \emph{aggregate over the array
range} \(P[\tin(v)\,..\,\tout(v)]\).
\end{definition}

\begin{theorem}{Subtree Mo is one flattened range sweep}{c21-subtree-range}
Given \(q\) offline queries of the form "statistic of subtree\((v_i)\)" where
the statistic supports \(O(1)\) element insert/delete but has no closed prefix
form, the block-order sweep of \secref{c21-moengine} over the ranges
\([\tin(v_i),\tout(v_i)]\) on array \(P\) answers all \(q\) in
\(\bigO{(n+q)\sqrt n}\).
\end{theorem}
\begin{proofbox}
The DFS visits every node of subtree\((v)\) between entering \(v\) and leaving
its region, and visits nothing else in between -- so subtree\((v)\) occupies
exactly the index span \([\tin(v),\tout(v)]\) of \(P\).  Under that bijection
node \(u\) is the single array position \(\tin(u)\): \texttt{add}/\texttt{rm}
touch one node of the current subtree-window per step.  The \emph{shape} of the
ranges is irrelevant to \thmref{c21-complexity} -- it bounded pointer movement
for any family of \(q\) ranges over an \(n\)-array.  Applying it verbatim gives
\(\bigO{(n+q)\sqrt n}\). \(\Box\)
\end{proofbox}

\begin{lstlisting}[caption={The flatten DFS plus the Mo-on-subtrees wiring: a query on subtree(v) is just the range tin[v] to tout[v] fed to the 27.1 block-order sweep. ASCII only, no underscores in identifiers}, label={lst:c21-flatten}]
int timer_, tin[MAXN], tout[MAXN], P[MAXN], dep[MAXN];
void dfs(int u, int p){
    tin[u] = ++timer_;  P[timer_] = u;
    for (int v : adj[u]) if (v != p){
        dep[v] = dep[u] + 1;
        dfs(v, u);
    }
    tout[u] = timer_;                    // subtree(u) = P[tin[u] .. tout[u]]
}
// Mo-on-subtrees: query i on subtree(v) becomes the range
//   l = tin[v],  r = tout[v]
// then sort by (l/B, r) and sweep with the add()/rm() of listing lst:c21-mo
// applied to P[pos] (the node at that preorder position).
\end{lstlisting}

\begin{tracebox}[Running example: subtree\((2)\) of \(R\) is one line of the flatten]
On \(R\) (\(\figref{c21-tree}\)) the DFS gives \(P=[1,2,4,8,5,3,6,7]\),
\(\tin(v)=(1,2,3,4,5,6,7,8)\) and \(\tout(v)=(8,5,6,4,5,7,8,4)\) for nodes
\(v=1..8\); in particular \(\tout(2)=5\).  subtree\((2)=\{2,4,8,5\}\) and indeed
\(P[2..5]=[2,4,8,5]\): a \emph{contiguous} range, though the four nodes are
scattered through the DFS.  A distinct-colour Mo would set
\(L,R\) targets inside \([2,5]\); the window's \texttt{add} of position \(i\)
inserts colour \(\mathit{val}(P[i])\): with \(\mathit{val}=(2,1,3,1,2,3,1,2)\)
from the opener, positions \(2,3,4,5\) hold nodes \(2,4,8,5\) with colours
\(1,3,2,2\) -- three distinct.  Nothing about the sweep changed; only the array
did. \(\Box\)
\end{tracebox}

\begin{insightbox}[The two worlds, decided by one question]
Ask: \emph{"can I get the answer on a range from two precomputed prefixes?"}
If yes (sums, xors, signed parity-corrected sums), a \(\bigO{\log n}\) BIT/
segment tree over the flatten crushes the problem and Mo is pure overhead.  If
no (distinct count, \(\sum \mathrm{freq}^2\), "how many values appear \(\ge l\)
times"), prefixes do not exist, but one-element \texttt{add}/\texttt{rm} do --
and that asymmetry is exactly the block-order sweep's salary.  CF~383C --
Ch.~18's solved case (\chref{c14}) -- is the study in spotting that a
\emph{scary-looking} update (sign flips with depth) lives in the first world.
CF~463E, solved below, lives in a \emph{third}: no inverse exists, yet no
reordering is needed either, because every queried root path is the current
recursion stack -- a nested window that a plain DFS sweeps for free.
\end{insightbox}

\begin{complexitybox}[Subtree Mo]
Flatten: \(\bigO{n}\) one DFS.  Per query: form its range in \(O(1)\).  Sweep:
\(\bigO{(n+q)\sqrt n}\) (\(\thmref{c21-subtree-range}\)).  Contrast: an
invertible statistic on the same flatten costs \(\bigO{(n+q)\log n}\) with a
difference BIT -- so \emph{never pay \(\sqrt n\) when a prefix exists}.  The
canonical non-invertible drill is CF~375D ("Tree and Queries": distinct colours
in subtrees, answered here by subtree ranges, in Ch.~17 by DSU-on-tree -- the
URL and full treatment live in \chref{c28}).
\end{complexitybox}

% ---------------------------------------------------------------
\subsection{Problem Lab: CF 463E Caisa and Tree -- the root-path stack and rare rebuilds}
% ---------------------------------------------------------------

\begin{problemstatement}[Caisa and Tree (statement, constraints and sample verified against the live CF page)]
A tree on \(n\le 100\,000\) nodes rooted at \(1\); node \(i\) has value
\(a_i\le 2\cdot 10^6\).  Process \(q\le 100\,000\) queries in order, two kinds:
\begin{itemize}[nosep]
  \item \texttt{1 v} -- write the root-to-\(v\) path as
        \(u_1,\dots,u_k\) (\(u_1=1\), \(u_k=v\)); output the vertex \(u_i\) with
        \(i<k\) and \(\gcd(a_{u_i}, a_v) > 1\) having the \emph{maximum} such
        \(i\); output \(-1\) if none exists;
  \item \texttt{2 v w} -- change \(a_v\) to \(w\).
\end{itemize}
The promise that shapes everything: \emph{at most \(50\) queries are of type 2}.

\textbf{Sample.} \(n{=}4,\ q{=}6\), \(a=[10,8,4,3]\), chain edges
\(1\!-\!2,\ 2\!-\!3,\ 3\!-\!4\); queries \(\texttt{1 1},\ \texttt{1 2},\
\texttt{1 3},\ \texttt{1 4},\ \texttt{2 1 9},\ \texttt{1 4}\) \(\to\)
\texttt{-1 1 2 -1 1}.
\end{problemstatement}

\begin{attemptgate}[Attempt it yourself first -- \lv{4}, give yourself 25 minutes]
\(\gcd(a_u,a_v)>1\) means "share \emph{some} prime".  You want the deepest such
ancestor for many \(v\), and the values barely change.  Can you avoid both
climbing per query and a lifting table?
\peek{Two numbers share a prime iff some prime \(p\mid\) both.  Factor \(a_v\)
into \emph{distinct} primes -- at most seven below \(2\cdot 10^6\).}
\nudge{For every prime \(p\) keep \(\mathit{last}[p]\): the deepest node on the
\emph{current root path} divisible by \(p\).  A DFS enters and leaves nodes, so
\(\mathit{last}\) is push-then-undo -- the \(O(1)\) \texttt{add}/\texttt{rm} of
27.1 again, this time riding the recursion stack instead of two pointers.}
\fullhint{Read \(v\)'s query \emph{before} pushing \(v\)'s primes -- the target
must be a strict ancestor (\(i<k\)).  Between two type-2 queries all values are
static, so cut the stream into \(\le 51\) blocks, run one DFS per block, apply
the single update, repeat.  Sieve smallest-prime-factors once up to
\(2\cdot 10^6\).  Cost: \(51\cdot n\) enters with \(\le 7\) table probes each
-- inside the 10-second limit with room to spare.  No Mo, no lifting: the
rebuild schedule \emph{is} the offline trade.}
\end{attemptgate}

\solstep{Understanding}
Two subproblems, each with a cheap answer.  (i) \emph{The statistic}:
"\(\gcd(a_u,a_v)>1\)" is existential over primes -- the deepest ancestor
divisible by \(p\), for some prime \(p\) of \(a_v\).  It is \emph{not} monotone
along the path (a grandparent may share a prime the parent does not), so binary
lifting has nothing to bite on -- that is ladder rung 2's WA.  A per-prime
\emph{last-seen} table answers it with at most seven probes.  (ii) \emph{The
schedule}: interleaved updates normally kill any offline sweep -- except there
are at most \(50\), so the query stream splits into \(\le 51\) static blocks and
a full re-DFS per block costs almost nothing compared to any sqrt-decomposition.

\solstep{Key Observations}
\begin{observation}
Every \(a\le 2\cdot 10^6\) has at most \(7\) distinct prime factors
(\(2{\cdot}3{\cdot}5{\cdot}7{\cdot}11{\cdot}13{\cdot}17 = 510510\), and
multiplying by the eighth prime exceeds \(2\cdot 10^6\)), so a factor list is a
constant-size key: \(O(1)\) probes and \(O(1)\) undos per node, always.
\end{observation}
\begin{observation}
During a DFS from the root the call stack \emph{is} the current root path:
writing \(\mathit{last}[p]=u\) on entry to \(u\) and restoring the saved value
on exit maintains "deepest ancestor divisible by \(p\)" for \emph{every} prime
at once.  Attach each type-1 query to its node and answer it at entry -- before
the node's own primes are pushed, so the match stays a strict ancestor.
\end{observation}

\solstep{The Approach Ladder}
\begin{center}\renewcommand{\arraystretch}{1.25}\setlength{\tabcolsep}{3pt}
\begin{tabular}{@{}c p{0.26\linewidth} p{0.30\linewidth} l c@{}}
\toprule
\textbf{\#} & \textbf{Approach} & \textbf{Idea} & \textbf{Time} & \textbf{Verdict}\\
\midrule
1 & \lv{2}\ climb the path per query & test \(\gcd(a_u,a_v)>1\) ancestor by ancestor & \(\bigO{qn}\) chain & \tle \\
2 & \lv{3}\ binary lifting on the path & lift to "deepest ancestor sharing a prime" & \(\bigO{q\log n}\) hope & \textbf{WA} \\
3 & \lv{4}\ \textbf{DFS stack + \(\mathit{last}[p]\) table} & one re-DFS per static block, \(O(1)\) push/undo & \(\bigO{51\,n\cdot 7+q}\) & \ac \\
\bottomrule
\end{tabular}\end{center}
Rung~2 is the trap the statement sets on purpose: binary lifting needs a
\emph{monotone} predicate, and "shares a prime with \(v\)" is not one -- a
grandparent can match where the parent fails, so the lift overshoots or
undershoots and no fix-up patch makes it right.  Rung~1 is honest but quadratic:
on the chain with \(n=q=10^5\) it is \(10^{10}\) gcds.  Rung~3 collects the two
gifts the statement hands you -- the queried window is always the current root
path (free to carry on the recursion stack) and updates number \(\le 50\) (so
the static-block rebuild is just \(51\) DFSes).

\solstep{Solution}
Sieve smallest-prime-factors \(\mathit{spf}[\cdot]\) up to \(2\cdot10^6\) once.
Cut the query stream into its \(\le 51\) maximal runs of type-1 queries
separated by type-2 updates.  For each run: attach every query to its vertex,
then DFS from \(1\) maintaining \(\mathit{last}[p]\) -- the deepest node on the
current stack divisible by prime \(p\).  On entry to \(u\): factor \(a_u\) into
\emph{distinct} primes; answer each query attached to \(u\) as the deepest
\(\mathit{last}[p]\) over those primes, \emph{read before} \(u\)'s own primes
are pushed (the match must be a strict ancestor); then save and overwrite
\(\mathit{last}[p]\leftarrow u\), recurse, and restore the saved values on
exit.  Print the block's answers in original order, apply the one update, and
start the next DFS.

\solstep{Implementation}
\begin{lstlisting}[caption={CF 463E: one DFS per static block. The recursion stack carries the root path; lastp[p] is the deepest stack node divisible by prime p, saved and undone around each node; queries are read before the push. ASCII only}, label={lst:c21-caisa}]
#include <bits/stdc++.h>
using namespace std;
const int MAXN=100005, MAXA=2000005;
int n,q,a[MAXN],spf[MAXA],dep[MAXN],ans[MAXN];
int typ[MAXN],qv[MAXN],qw[MAXN],lastp[MAXA];
vector<int> adj[MAXN],ask[MAXN];

void build_spf(){
  for(int i=2;i<MAXA;i++){
    if(!spf[i]){
      spf[i]=i;
      if((long long)i*i<MAXA)
        for(long long j=(long long)i*i;j<MAXA;j+=i)
          if(!spf[j]) spf[j]=i;
    }
  }
}
void factor(int x, vector<int>& out){      // distinct primes only
  out.clear();
  while(x>1){ int p=spf[x]; out.push_back(p); while(x%p==0) x/=p; }
}
void dfs(int u,int p){
  vector<int> f; factor(a[u],f);
  int best=0;
  for(int pr:f)                              // deepest over primes of a[u]
    if(lastp[pr] && (!best || dep[lastp[pr]]>dep[best])) best=lastp[pr];
  for(int id:ask[u]) ans[id] = best ? best : -1;   // BEFORE pushing: i<k
  vector<pair<int,int> > saved;
  for(int pr:f){ saved.push_back(make_pair(pr,lastp[pr])); lastp[pr]=u; }
  for(int v:adj[u]) if(v!=p){ dep[v]=dep[u]+1; dfs(v,u); }
  for(int k=0;k<(int)saved.size();k++)       // undo on exit
    lastp[saved[k].first]=saved[k].second;
}
int main(){
  build_spf();
  scanf("%d%d",&n,&q);
  for(int i=1;i<=n;i++) scanf("%d",&a[i]);
  for(int i=1;i<n;i++){ int u,v; scanf("%d%d",&u,&v);
                        adj[u].push_back(v); adj[v].push_back(u); }
  for(int i=0;i<q;i++){
    scanf("%d%d",&typ[i],&qv[i]);
    if(typ[i]==2) scanf("%d",&qw[i]);
  }
  int id=0;
  while(id<q){
    int start=id;
    for(int u=1;u<=n;u++) ask[u].clear();    // fresh attachment per block
    while(id<q && typ[id]==1){ ask[qv[id]].push_back(id); id++; }
    dep[1]=0; dfs(1,0);                      // whole static block, one DFS
    for(int j=start;j<id;j++) printf("%d\n",ans[j]);
    if(id<q){ a[qv[id]]=qw[id]; id++; }      // the single update between blocks
  }
  return 0;
}
\end{lstlisting}

\solstep{Trace}
\begin{tracebox}[Replaying the official sample to \(-1,1,2,-1,1\)]
Chain \(1\!-\!2\!-\!3\!-\!4\), \(a=[10,8,4,3]\).  Block~1 attaches queries
1--4 to nodes 1--4; \(\mathit{last}[\cdot]\) starts empty. \quad
\textbf{enter~1} (\(10=2{\cdot}5\)): probe \(\mathit{last}[2],\mathit{last}[5]\)
-- both empty \(\Rightarrow\) query~1 \(=\boxed{-1}\); push
\(\mathit{last}[2]=\mathit{last}[5]=1\). \quad
\textbf{enter~2} (\(8=2^3\to\{2\}\)): \(\mathit{last}[2]=1\Rightarrow\)
query~2 \(=\boxed{1}\); push \(\mathit{last}[2]=2\). \quad
\textbf{enter~3} (\(4\to\{2\}\)): \(\mathit{last}[2]=2\Rightarrow\) query~3
\(=\boxed{2}\); push \(3\). \quad
\textbf{enter~4} (\(3\to\{3\}\)): \(\mathit{last}[3]\) empty \(\Rightarrow\)
query~4 \(=\boxed{-1}\).  Every exit restores \(\mathit{last}\) to empty. \quad
Apply \texttt{2 1 9}: \(a_1=9=3^2\).  Block~2 attaches query~5 to node~4.
\textbf{enter~1}: push \(\mathit{last}[3]=1\); nodes 2,3 push
\(\mathit{last}[2]=2,3\); \textbf{enter~4} (\(\{3\}\)):
\(\mathit{last}[3]=1\Rightarrow\) query~5 \(=\boxed{1}\) -- the updated
\(a_1=9\) now shares the prime 3 with \(a_4=3\).  Output
\texttt{-1 1 2 -1 1} -- match. \(\Box\)
\end{tracebox}

\begin{warningbox}[Three ways CF 463E goes wrong]
\begin{enumerate}[nosep,leftmargin=1.6em]
  \item \textbf{Answering after the push.} If \(u\)'s own primes are already
        in the table when its queries are read, \(v\) matches \emph{itself} --
        the statement demands \(i<k\), so every non-\(-1\) answer degenerates
        to \(v\).
  \item \textbf{Factoring with multiplicity.} Pushing prime \(2\) three times
        for \(a_u=8\) makes the undo restore the \emph{middle} copy: after
        leaving \(u\), \(\mathit{last}[2]\) still points at \(u\), and every
        later probe in the DFS is contaminated by a node off the current path.
  \item \textbf{Stale block state.} Attachments (\(\mathit{ask}\)) not cleared
        between blocks fire twice; an update applied in the wrong order (after
        the DFS that should have seen it, or never applied) shifts whole
        answer runs by one version.
\end{enumerate}
\end{warningbox}

\begin{hardeningbox}
\begin{itemize}[nosep,leftmargin=1.4em]
  \item \textbf{The stack window is 27.1 in disguise.} Push-on-entry /
        undo-on-exit is \texttt{add}/\texttt{rm} over a \emph{nested} window.
        You cannot reorder DFS visits the way Mo reorders ranges, so the
        primitive pays exactly when the queried sets are root paths.  The
        Summit (\(\probref{c21-summit}\)) is the same skeleton carrying a full
        frequency table instead of a last-seen table.
  \item \textbf{Price the rebuild honestly.} \((U{+}1)\) full DFSes is the
        trade when updates number \(U\le 50\); it degrades gracefully to
        \(U=\bigO{\sqrt n}\), but \(U=\bigO{n}\) kills it -- that world is Mo
        with modifications or an online structure, not this chapter's door.
  \item \textbf{Table sizing.} \(\mathit{last}[\cdot]\) is indexed by
        \emph{prime value}, not by node: it needs \(2\cdot10^6\) slots, while
        everything else is \(\bigO{n}\).  Sizing it by \(n\) is an
        out-of-bounds write that samples will not catch.
  \item \textbf{Depth bookkeeping.} Comparing "deepest" across primes needs
        \(\depth\); fill it on the same DFS.  A \(10^5\)-deep chain is fine on
        Codeforces' stack, but the explicit-stack version is the
        belt-and-braces port (same caution as Tarjan, \(\secref{c21-tarjan}\)).
\end{itemize}
\end{hardeningbox}

% ---------------------------------------------------------------
\subsection{The Mo-on-subtrees recipe -- and the two doors that beat it}
% ---------------------------------------------------------------

\begin{recognitionbox}[Subtree queries: which tool, exactly]
\begin{itemize}[nosep,leftmargin=1.4em]
  \item \textbf{Distinct colours / freq-shape statistic over subtree\((v)\),
        all queries known upfront.}
        \reckey{run \secref{c21-moengine}'s sweep on ranges
        \([\tin(v),\tout(v)]\) of \(P\); CF 375D is the drill, treated in
        \chref{c28}.}
        \recuse{\item window \(=\) current subtree span; \texttt{add} \(=\)
        insert one node's colour.}
        \recnotuse{\item never re-DFS per query.}
  \item \textbf{Signed/invertible subtree update, point or range read.}
        \reckey{flatten + difference BIT, \(\bigO{\log n}\) per operation -- the
        solved template is Ch.~18's CF 383C (\(\chref{c14}\)).}
        \recuse{\item the parity factor \((-1)^{\depth}\) when signs alternate.}
        \recnotuse{\item Mo here is \(O(\sqrt n)\) tax on a problem with change.}
  \item \textbf{Online updates mixed with subtree aggregates, subtrees only.}
        \reckey{Euler + segment tree; the path version needs the doubled tour
        (\(\secref{c21-path}\)) or HLD (\chref{c16}).}
        \recuse{\item the segment tree inherits the same \([\tin,\tout]\)
        ranges.}
        \recnotuse{\item block order cannot answer query \(i\) after an update
        at query \(i' > i\).}
  \item \textbf{Root-path statistic, strict-ancestor flavour, rare updates.}
        \reckey{DFS-stack window + per-key push/undo table, one rebuild per
        static block (CF 463E, \(\probref{c21-caisa}\)).}
        \recuse{\item read the answer \emph{before} pushing the node's own keys.}
        \recnotuse{\item no binary lifting -- an existential predicate over
        primes is not monotone; no Mo -- a DFS already visits every root path
        in order.}
\end{itemize}
\end{recognitionbox}

\begin{notebox}[The one-line transfer principle]
Every array-Mo solution from 27.1 becomes a subtree-Mo solution by the
substitution \(l\to\tin(v),\ r\to\tout(v)\) -- the proof of
\thmref{c21-subtree-range} is that the sweep never inspects the ranges'
\emph{meaning}.  The next section performs the harder transfer, to
\emph{paths}, where no flatten makes the target contiguous and a
\term{toggle} must be invented.
\end{notebox}

\begin{recallbox}[27.2 in one breath]
\emph{Subtree\((v)\) is the range \([\tin(v),\tout(v)]\) of the preorder array
\(P\) (\(\defref{c21-flatten}\)) -- non-invertible subtree statistics ride the
27.1 block-order sweep unchanged (\(\thmref{c21-subtree-range}\),
\(\coderef{c21-flatten}\)); invertible ones skip Mo (Ch.~18's CF 383C,
\(\chref{c14}\)); and CF 463E opens the third door: the queried set is the
\emph{current root path}, so the recursion stack carries the window -- a
per-prime \(\mathit{last}[p]\) table with \(O(1)\) push/undo, answers read
\emph{before} the push, one re-DFS per block of the \(\le 50\) updates
(\(\probref{c21-caisa}\), \(\coderef{c21-caisa}\), sample replayed to
\(-1,1,2,-1,1\)).}
\selftest{
  \item Which world does "distinct colours in subtree\((v)\), offline" live in, and what answers it? \emph{(No inverse -- Mo's sweep over \([\tin(v),\tout(v)]\) ranges, or DSU-on-tree; \thmref{c21-subtree-range}. \secref{c21-boundary} contrasts them.)}
  \item Why does CF 463E need neither Mo nor binary lifting? \emph{(The DFS stack already visits every root path -- there is no reordering to gain; and "shares a prime with \(v\)" is non-monotone, so lifting has no valid predicate.)}
  \item Name the two ordering facts that make the push/undo table correct. \emph{(Read each node's queries before pushing its own primes -- strict ancestor \(i<k\); factor into distinct primes and undo to the saved values on exit -- one live entry per prime.)}
  \item What promise buys the \(51\) rebuilds, and what would kill the approach? \emph{(\(\le 50\) type-2 queries; at \(U=\bigO{n}\) updates the rebuild is \(\bigO{nq}\) -- that world needs Mo-with-modifications or an online structure.)}
}
\end{recallbox}
